\subsubsection{Related Work}
In our earlier work \cite{DTS-2025} we developed an algorithm for
solving the problem
of realisability of sets of (sequential) scenarios as distributed scenarios.
Given a set of scenarios, our algorithm determines whether
the set is realisable by a \emph{single} DTS, i.e.,
whether there is a DTS whose semantics is identical
to the union of the behaviours described by the members of
the set. If so, the algorithm produces such a DTS.
Despite some similarities, the problem that we tackle in the current
paper is fundamentally different.

%----------------------
The comparison of our sequential timed scenarios
with simple temporal networks (STNs) \cite{DECHTER1991}, as well as
with message sequence charts (MSCs) \cite{Harel2003,MSCs}, is presented
elsewhere \cite{DTS-2025}.
Similarities between
%our distance relations and timed partial
%orders (TPOs) of Watanabe et al. \cite{watanabe2023}, as well as
%between
the distributed behaviours of distributed timed scenarios and the
timed traces of Balaguer et al. \cite{hal-2012} are also discussed
in our earlier work \cite{DTS-2025}.

%The comparison of our stable distance tables with Difference Bounds
%Matrices is presented elsewhere \cite{neda-2017,neda-2018-2}.

The problem of realisability has been investigated in a variety of settings,
%of global models
%has been investigated
e.g., for message sequence
charts \cite{Alur-2003}, global session
types \cite{BEJLERI-2009,Castagna-2011,Honda-2016,Hans-2016},
pomsets \cite{GUANCIALE201969},
global labelled transition systems \cite{Beek-2023},
%more recently for
choreography languages \cite{Franco-2022} and team
automata \cite{Beek-2024}. Unlike in all
these works the realisability of our sets of sequential scenarios as
scenario expressions is
determined by precise time analysis enabled by our distance
relations.
%The distance relations enable us to reason about the underlying partail
%order, as well as the time constraints between various events in
%scenarios, both of which are essential for realisability of
That is, in our setting we are able to directly reason
about realisability of models that include explicit time constraints
and produce realisations that involve time constraints.
